% ID: SEL-D0039-OA-E01
% ESTADO: validada. S040, cotejada con examen y criterios.
Sean $x,y$ las tarrinas del primer y segundo tipo. En gramos,
los límites de chocolate y straciatella son
$100x+150y\le8000$ y $200x+150y\le10000$.
Tras simplificar y añadir los barquillos:
\[
2x+3y\le160,\qquad 4x+3y\le200,\qquad x+2y\le100,
\qquad x,y\ge0.
\]
Las tres fronteras coinciden en $(20,40)$. En los ejes los
cortes factibles son $(50,0)$ y $(0,50)$. El cuarto vértice
es el origen.
\begin{center}
\setlength{\unitlength}{7mm}
\begin{picture}(10,8)(-2,-1)
\qbezier(0,0)(2.8,0)(5.6,0)\put(5.7,0){$x$}
\qbezier(0,0)(0,2.8)(0,5.6)\put(0,5.7){$y$}
\qbezier(0,5)(1,4.5)(2,4)
\qbezier(2,4)(3.5,2)(5,0)
\put(-0.15,5){\makebox(0,0)[r]{$(0,50)$}}
\put(2.1,4.2){$(20,40)$}
\put(5,-0.3){\makebox(0,0)[t]{$(50,0)$}}
\end{picture}
\end{center}
El total $x+y$ toma los valores $0,50,60,50$ en
$(0,0)$, $(50,0)$, $(20,40)$ y $(0,50)$.
Por linealidad, el mejor vértice da el máximo. Sus coordenadas
son enteras: se preparan \textbf{20 del primer tipo y 40 del
segundo}, un total de $\boxed{60}$ tarrinas.
